\section*{第 \ref{ch:g12:arith} 章——算术}

\begin{solution}{exo:g12:arith:1}
$17 \times 119 = 2023$，故 $2026 = 17 \times 119 + 3$：商 $119$，
余数 $3$。对 $-2026$：$-2026 = 17\times(-120) + 14$（事实上
$17 \times 120 = 2040$，且 $2040 - 2026 = 14$）：商 $-120$，余数
$14$（余数须属于 $\intco{0}{17}$，故它\emph{不是} $-3$）。
\end{solution}

\begin{solution}{exo:g12:arith:2}
模 $10$：$7^2 = 49 \equiv 9 \equiv -1$。因此
$7^{100} = \left(7^2\right)^{50} \equiv (-1)^{50} = 1 \pmod{10}$：$7^{100}$ 的末位
数字为 $1$。
\end{solution}

\begin{solution}{exo:g12:arith:3}
欧几里得算法：$1071 = 2\times462 + 147$；$462 = 3\times147 + 21$；
$147 = 7\times21 + 0$。故 $\gcd = 21$。

回代：$21 = 462 - 3\times147 = 462 - 3(1071 - 2\times462)
= 7\times462 - 3\times1071$。于是 $u = -3$，$v = 7$：
$1071\times(-3) + 462\times7 = 21$。
\end{solution}

\begin{solution}{exo:g12:arith:4}
每个\omterm{def:g10:numbers:sets}{整数}都满足 $\equiv 0, 1, 2$ 或 $3 \pmod 4$，并计算平方：
$0^2 \equiv 0$，$1^2 \equiv 1$，$2^2 = 4 \equiv 0$，$3^2 = 9 \equiv 1$。故
$n^2 \equiv 0$ 或 $1 \pmod 4$。于是两个平方数之和\omterm{def:g12:arith:congruence}{同余}于
$0 + 0$、$0 + 1$ 或 $1 + 1$，\emph{即}为 $0$、$1$ 或 $2 \pmod 4$——
决不与 $3$ \omterm{def:g12:arith:congruence}{同余}。
\end{solution}

\begin{solution}{exo:g12:arith:5}
被 $2$ \omterm{def:g12:arith:divides}{整除}：$n$ 与 $n + 1$ 中必有一个为偶数。被
$3$ \omterm{def:g12:arith:divides}{整除}：若 $n \equiv 0$，则 $3 \mid n$；若 $n \equiv 1 \pmod 3$，则
$2n + 1 \equiv 3 \equiv 0$；若 $n \equiv 2$，则 $n + 1 \equiv 0$。于是各情形下
$3$ \omterm{def:g12:arith:divides}{整除}该乘积。由于 $\gcd(2,3) = 1$，
由\cref{cor:g12:arith:coprimeprod}可得 $6 \mid n(n+1)(2n+1)$。
（这也重新证明了 $\frac{n(n+1)(2n+1)}{6}$，即
\cref{exo:g12:seq:1}中的平方和，是\omterm{def:g10:numbers:sets}{整数}。）
\end{solution}

\begin{solution}{exo:g12:arith:6}
求 $5$ 模 $11$ 的逆元：试算（或用贝祖等式）
$5 \times 9 = 45 = 44 + 1 \equiv 1 \pmod{11}$。将该\omterm{def:g12:arith:congruence}{同余}式两边乘以
$9$：
\[
x \equiv 9 \times 3 = 27 \equiv 5 \pmod{11}.
\]
所有解为\omterm{def:g10:numbers:sets}{整数} $x = 5 + 11k$，$k \in \Z$。（验算：
$5\times5 = 25 \equiv 3 \pmod{11}$。）
\end{solution}

\begin{solution}{exo:g12:arith:7}
$\gcd(17, 40) = 1$，故解存在。欧几里得算法：
$40 = 2\times17 + 6$；$17 = 2\times6 + 5$；$6 = 5 + 1$。回代：
$1 = 6 - 5 = 6 - (17 - 2\times6) = 3\times6 - 17
= 3(40 - 2\times17) - 17 = 3\times40 - 7\times17$。因此
$17\times(-7) - 40\times(-3) = 1$，故特解
为 $(x_0, y_0) = (-7, -3)$。

\omterm{def:g10:algebra:equation}{方程} $17x - 40y = 1$ 的通解：将特解关系式相减，
$17(x + 7) = 40(y + 3)$；由于 $\gcd(17, 40) = 1$，由高斯引理可得
$40 \mid x + 7$，故 $x = -7 + 40k$，进而 $y = -3 + 17k$，$k \in \Z$
（经验证均成立）。

对于 $17x - 40y = 6$，将特解乘以 $6$：
$(x_1, y_1) = (-42, -18)$，同理可得
\[
x = -42 + 40k, \qquad y = -18 + 17k, \qquad k \in \Z .
\]
（例如 $k = 2$：$x = 38$，$y = 16$；事实上 $17\times38 - 40\times16
= 646 - 640 = 6$。）
\end{solution}

\begin{solution}{exo:g12:arith:8}
设 $\sqrt2 = \frac ab$，其中 $a, b \in \N^*$，则 $a^2 = 2b^2$。在
平方数的质因数分解中，每个指数都是偶数，故
$2$ 在 $a^2$ 中的指数为偶数，而在 $2b^2$ 中为奇数（比某个偶数多一）。同一\omterm{def:g10:numbers:sets}{整数}的两个
质因数分解中 $2$ 的指数不同，
与\cref{thm:g12:arith:fta}的唯一性相矛盾。因此
不存在这样的分数：$\sqrt2 \notin \Q$。
\end{solution}

\begin{solution}{exo:g12:arith:9}
\emph{1.} 由 $k\binom pk = p \binom{p-1}{k-1}$，$p$ \omterm{def:g12:arith:divides}{整除}
$k\binom pk$。对 $1 \leq k \leq p-1$，$p \nmid k$，且 $p$ 为\omterm{def:g12:arith:prime}{质数}，于是
$\gcd(p, k) = 1$，由高斯引理可得 $p \mid \binom pk$。

\emph{2.} 对 $a$ 作归纳。当 $a = 0$ 时，$0^p \equiv 0$。假设
$a^p \equiv a \pmod p$。由二项式定理，
\[
(a+1)^p = \sum_{k=0}^{p} \binom pk a^k
\equiv a^p + 1 \pmod p,
\]
由第 1 小题，中间各项模 $p$ 均消失。由归纳
假设，$(a+1)^p \equiv a + 1 \pmod p$。这就证明了
$a^p \equiv a$ 对所有 $a \in \N$ 成立，而 $a < 0$ 的情形由
$a \equiv a + kp$（取适当的正代表元）即得。
\end{solution}

\begin{solution}{exo:g12:arith:10}
$n \equiv 2 \pmod 3$ 且 $n \equiv 3 \pmod 5$：记 $n = 2 + 3s$，则
$2 + 3s \equiv 3 \pmod 5$，\emph{即} $3s \equiv 1 \pmod 5$。
$3$ 模 $5$ 的逆元为 $2$（$3\times2 = 6 \equiv 1$），故 $s \equiv 2 \pmod 5$，
记 $s = 2 + 5t$，于是 $n = 8 + 15t$。前两个条件即
$n \equiv 8 \pmod{15}$。

再加上 $n \equiv 2 \pmod 7$：$8 + 15t \equiv 2 \pmod 7$，且
$15 \equiv 1 \pmod 7$，故 $t \equiv -6 \equiv 1 \pmod 7$，记
$t = 1 + 7u$。因此 $n = 23 + 105u$：
\[
n \equiv 23 \pmod{105}.
\]
（验算：$23 = 3\times7 + 2 = 5\times4 + 3 = 7\times3 + 2$。）
\end{solution}

\begin{solution}{pb:g12:arith:1}
\textbf{1.} $2026 = 289 \times 7 + 3$：$2026 \equiv 3
\pmod 7$。$7$ 的各次幂模 $10$ 依次为 $7, 9, 3, 1$，循环节长为 $4$；
$100 \equiv 0 \pmod 4$：$7^{100}$ 的末位
数字为 $1$。

\textbf{2.} $5^2 = 25 \equiv -1 \pmod{13}$，故
$5^{116} = \left(5^2\right)^{58} \equiv (-1)^{58} = 1$，且
$5^{117} \equiv 5 \pmod{13}$。

\textbf{3.} $3$ 模 $7$ 的逆元为 $5$
（$15 \equiv 1$）：$x \equiv 5 \times 5 = 25 \equiv 4 \pmod 7$。

\textbf{4.} $97 = 2 \times 35 + 27$；$35 = 27 + 8$；
$27 = 3 \times 8 + 3$；$8 = 2 \times 3 + 2$；$3 = 2 + 1$。
回代得 $1 = 97 \times 13 + 35 \times (-36)$。于是
$35 \times (-36) \equiv 1 \pmod{97}$，故 $35$ 的逆元为
$-36 \equiv 61 \pmod{97}$。

\textbf{5.} $a$ 模 $n$ 可逆，当且仅当
$\gcd(a, n) = 1$：贝祖等式给出 $au + nv = 1$，即\
$au \equiv 1$；反之，若存在逆元，则\omterm{def:g12:arith:gcd}{最大公因数}必\omterm{def:g12:arith:divides}{整除}
$1$。

\textbf{6.} $0{\cdot}10 + 3{\cdot}9 + 0{\cdot}8 + 6{\cdot}7 +
4{\cdot}6 + 0{\cdot}5 + 6{\cdot}4 + 1{\cdot}3 + 5{\cdot}2 +
2{\cdot}1 = 132 = 12 \times 11 \equiv 0 \pmod{11}$：通过校验。

\textbf{7.} 和的改变量为 $wd$，其中
$1 \leq w \leq 10$ 且 $1 \leq \abs d \leq 9$：由于 $11$ 是
\omterm{def:g12:arith:prime}{质数}，且不\omterm{def:g12:arith:divides}{整除}任一因子，故其不能\omterm{def:g12:arith:divides}{整除}乘积
（见\cref{thm:g12:arith:gauss} / \cref{prop:g12:arith:primedivides}）：
改变后的和永远不再 $\equiv 0$：一位数字的
错误都会触发报警。

\textbf{8.} 交换相邻数字 $a, b$（权重分别为
$w + 1, w$）时，和的改变量为
$(w+1)b + wa - (w+1)a - wb = b - a \not\equiv 0$，
当 $a \neq b$ 时即被检出。其奥秘在于\emph{\omterm{def:g12:arith:prime}{质数}性}
（$11$ 为\omterm{def:g12:arith:prime}{质数}）：模 $10$ 时，像 $5 \times 2$ 这样的乘积
竟为零（两因子皆不为零），因此权重为 $5$ 的 $\pm 2$ 型错误
（或某种不幸的换位）可能漏检。

\textbf{9.} 十二位数字的加权和为 $119$；
校验位须使之成为 $10$ 的倍数：校验位为 $1$（完整
代码 $978\,2940199\,051$）。EAN 漏检的相邻换位
满足 $2(a - b) \equiv 0 \pmod{10}$，即\
$\abs{a - b} = 5$：例如交换 $2$ 与 $7$，会
被漏检——这是选用便捷模数 $10$ 付出的代价。

\textbf{10.} 自右向左每隔一位的数字加倍，
并把两位数的数字相加（$16 \to 7$ 等），总和为 $80 \equiv 0
\pmod{10}$：检验卡通过校验。

\textbf{11.} 模数为\omterm{def:g12:arith:prime}{质数}时，每个权重均可逆，
故\emph{所有}单个错误与\emph{所有}相邻换位
均可被发现——ISBN 的优势；模 $10$ 的方案
保留便于人读的数字，但接受了一个小的盲区。

\textbf{12.} $2^{10} = 1024 = 3 \times 341 + 1 \equiv 1
\pmod{341}$，故 $2^{340} = \left(2^{10}\right)^{34} \equiv
1$。然而 $341 = 11 \times 31$ 为合数：它通过了
以 $2$ 为底的费马检验，却并非\omterm{def:g12:arith:prime}{质数}。寓意是：费马\omterm{def:g12:arith:congruence}{同余}式
是必要的而非充分的——\omterm{def:g12:arith:prime}{质数}判定
需要更精细的工具（它们在
大学分册中给出）。

\textbf{13.} $3d \equiv 1 \pmod{20}$：$d = 7$
（$21 = 20 + 1$）。

\textbf{14.} $c = 4^3 = 64 \equiv 31 \pmod{33}$。

\textbf{15.} $31 \equiv -2$：$(-2)^7 = -128$，且
$-128 + 4 \times 33 = 4$：密文解密得
$m = 4$。锁开了。

\textbf{16.} 模 $3$ 时，若 $3 \nmid m$，
$m^2 \equiv 1$（费马），故
$m^{21} = m \cdot \left(m^2\right)^{10} \equiv m$；若
$3 \mid m$，则两边均 $\equiv 0$。模 $11$ 时：
$m^{10} \equiv 1$ 或 $11 \mid m$，且
$m^{21} = m \cdot \left(m^{10}\right)^2 \equiv m$。$3$ 与
$11$ 均\omterm{def:g12:arith:divides}{整除} $m^{21} - m$，且二者\omterm{def:g12:arith:gcd}{互质}，故其乘积
$33$ 亦\omterm{def:g12:arith:divides}{整除}（高斯）：$m^{21} \equiv m \pmod{33}$。
指数 $ed = 21 = 1 + 20k$ 的设计使两个费马
指数（$2$ 与 $10$，均\omterm{def:g12:arith:divides}{整除} $20$）消失。

\textbf{17.} 两个 $300$ 位\omterm{def:g12:arith:prime}{质数}相乘，仅需
一微秒；而由乘积还原它们，则击败了每一种
已知算法及全世界的计算机——这把锁是
一条单行道。（我们的 $n = 33$ 只是玩具规模下的
双向可通的道路。）

\textbf{18.} 试算余数（或用贝祖等式构造）：
$n \equiv 8 \pmod{15}$：计数依次为 $8, 23, 38, 53, \dots$
模 $15$ 的唯一性：两解之差是
$3$ 与 $5$ 的公倍数，从而为 $15$ 的倍数（$3$ 与 $5$ \omterm{def:g12:arith:gcd}{互质}，高斯）。
率领 $1000$ 名士兵的将军，经三次快速排队宣布“$8$”
——这正是古老的点兵术。

\textbf{19.} 由 $10 \equiv 1 \pmod 9$ 得 $10^k \equiv 1$，故
$\sum d_k 10^k \equiv \sum d_k$：一个数与其数字和
模 $9$ \omterm{def:g12:arith:congruence}{同余}（亦模 $3$ \omterm{def:g12:arith:congruence}{同余}）。又
由 $10 \equiv -1 \pmod{11}$ 得
$\sum d_k 10^k \equiv \sum (-1)^k d_k$：即交替符号法则。
两条童年时学过的法则，各用一行即可。

\textbf{20.} \omterm{def:g12:arith:congruence}{同余}把余数变成了一门算术
（第一部分）；贝祖等式造出了解线性\omterm{def:g12:arith:congruence}{同余}式及 RSA
的 $d$ 所需的逆元（第 4、13 题）；费马小定理
开启并闭合了那把锁（第 15--16 题）；拼接
\omterm{def:g12:arith:gcd}{互质}模数统计了士兵数，并完成了证明
（第 16、18 题）。关于哈代的评判：他所知最纯粹的定理
如今守护着每一笔购买——纯粹，假以时日，便是
世上最具应用价值的东西。
\end{solution}
